%! Sinusoidal Functions — Unit 6, Lesson 6.6 — Group Activity
\documentclass[10pt]{article}
\usepackage{precalculus-article}
\usepackage{precalculus-boxes}
\newcommand{\TallMath}[1]{$\displaystyle #1\rule[-1.4em]{0pt}{3.2em}$}

\begin{document}

\pageheader{Unit 6, Lesson 6.6}{Group Activity}
\namepartnerperiod

\vspace{0.1in}

\begin{scenariobox}[Sinusoidal Functions in the Real World]{sky}
Sinusoidal functions appear everywhere in nature and technology: sound waves,
tides, seasonal temperature patterns, alternating electric current, and
the motion of a pendulum.  In each case, the key characteristics ---
amplitude, midline, period, and frequency --- tell us something meaningful
about the phenomenon.
\end{scenariobox}

\vspace{0.1in}

%=======================================================================
% TIER R
%=======================================================================
\begin{tcolorbox}[colback=white, colframe=black!40,
    title=\textbf{Tier R --- Remediate},
    fonttitle=\bfseries, arc=2mm, left=3mm, right=3mm, top=2mm, bottom=2mm]

Use the formulas:
\[
  \text{Amplitude} = \frac{\text{max} - \text{min}}{2}
  \qquad
  \text{Midline: } y = \frac{\text{max} + \text{min}}{2}
  \qquad
  \text{Frequency} = \frac{1}{\text{period}}
\]

\textbf{Function 1.} A sinusoidal function has a maximum value of $4$, a minimum
value of $-4$, and completes one full cycle from $\theta = 0$ to $\theta = 2\pi$.

\medskip
\begin{tabular}{@{}llll@{}}
  Amplitude $=$ \blank{2cm} &
  Midline $=$ \blank{2cm} &
  Period $=$ \blank{2cm} &
  Frequency $=$ \blank{2.5cm}
\end{tabular}

\vspace{0.15in}

\textbf{Function 2.} A sinusoidal function has a maximum value of $3$ and a
minimum value of $-1$.

\medskip
\begin{tabular}{@{}ll@{}}
  Amplitude $=$ \blank{2cm} &
  Midline $=$ \blank{2cm}
\end{tabular}

\vspace{0.1in}

Is the midline $y = 0$ here?  Why or why not?
\writelines{2}

\end{tcolorbox}

\vspace{0.1in}

%=======================================================================
% TIER A
%=======================================================================
\begin{tcolorbox}[colback=white, colframe=black!40,
    title=\textbf{Tier A --- At Level},
    fonttitle=\bfseries, arc=2mm, left=3mm, right=3mm, top=2mm, bottom=2mm]

Two sinusoidal functions are described by the tables below.

\medskip
\textbf{Function P:}

\renewcommand{\arraystretch}{1.4}
\begin{tabular}{|c|c|c|c|c|c|c|c|}
\hline
$\theta$ & $0$ & $\pi/2$ & $\pi$ & $3\pi/2$ & $2\pi$ & $5\pi/2$ & $3\pi$ \\
\hline
$P(\theta)$ & $3$ & $5$ & $3$ & $1$ & $3$ & $5$ & $3$ \\
\hline
\end{tabular}

\medskip
\textbf{Function Q:}

\begin{tabular}{|c|c|c|c|c|c|c|c|}
\hline
$\theta$ & $0$ & $\pi/4$ & $\pi/2$ & $3\pi/4$ & $\pi$ & $5\pi/4$ & $3\pi/2$ \\
\hline
$Q(\theta)$ & $0$ & $2$ & $0$ & $-2$ & $0$ & $2$ & $0$ \\
\hline
\end{tabular}

\medskip
For each function, determine the amplitude, midline, period, and frequency.

\medskip
\begin{tabular}{@{}lllll@{}}
  & \textbf{Amplitude} & \textbf{Midline} & \textbf{Period} & \textbf{Frequency} \\[4pt]
  Function $P$: & \blank{1.8cm} & \blank{1.8cm} & \blank{1.8cm} & \blank{2.2cm} \\[8pt]
  Function $Q$: & \blank{1.8cm} & \blank{1.8cm} & \blank{1.8cm} & \blank{2.2cm}
\end{tabular}

\vspace{0.12in}

Based on the table values, determine whether each function appears to be odd, even,
or neither.  Justify your answer using the definitions $f(-\theta) = f(\theta)$
(even) and $f(-\theta) = -f(\theta)$ (odd).

\medskip
Function $P$ is \blank{2cm} because: \writelines{2}

Function $Q$ is \blank{2cm} because: \writelines{2}

\end{tcolorbox}

\vspace{0.1in}

%=======================================================================
% TIER E
%=======================================================================
\begin{tcolorbox}[colback=white, colframe=black!40,
    title=\textbf{Tier E --- Extend},
    fonttitle=\bfseries, arc=2mm, left=3mm, right=3mm, top=2mm, bottom=2mm]

\textbf{Part 1 --- Reversing the Process.}
A sinusoidal function has amplitude $2$, midline $y = 1$, and period $\pi$.

\medskip
(a) What are the maximum and minimum values of this function?

\medskip
\textit{Hint:} Use the relationships
max $= $ midline $+$ amplitude and min $= $ midline $-$ amplitude.

\writelines{3}

\medskip
(b) What is the frequency of this function?

\writelines{2}

\vspace{0.1in}

\textbf{Part 2 --- Proving a Sine Identity.}
We know $\sin(-\theta) = -\sin\theta$ (sine is odd).

\medskip
Use the fact that on the unit circle, the point at angle $\pi - \theta$ has the
same $y$-coordinate as the point at angle $\theta$ to explain why
\[
  \sin(\pi - \theta) = \sin\theta.
\]

Write a complete explanation (at least two sentences).

\writelines{5}

\end{tcolorbox}

\end{document}
